\chapter{Communication Theory and System Model}

\section{Complex Baseband Representation}
The passband signal received by a digital receiver is commonly converted to
complex baseband form:
\begin{equation}
r[n] = I[n] + jQ[n],
\end{equation}
where $I[n]$ and $Q[n]$ are the in-phase and quadrature components. This
representation preserves amplitude and phase information while allowing
digital signal processing at a lower sampling rate.

For a transmitted symbol sequence $x[n]$, a simplified discrete-time channel
model is
\begin{equation}
y[n] = h[n]x[n]e^{j(2\pi \Delta f n+\phi)} + w[n],
\end{equation}
where $h[n]$ is the complex channel coefficient, $\Delta f$ is the
normalized frequency offset, $\phi$ is the carrier phase offset, and $w[n]$
is complex AWGN. In the Phase 2 experiment, $h[n]$ is constant over one
frame and changes independently between frames, which models block
Rayleigh fading.

\section{Digital Modulation}
Digital modulation maps groups of binary information bits to complex
constellation points. BPSK uses two phase states and carries one bit per
symbol. QPSK uses four phase states and carries two bits per symbol. 8-PSK
carries three bits per symbol, while normalized 16-QAM uses sixteen
amplitude-phase combinations and carries four bits per symbol.

The number of bits per symbol affects both spectral efficiency and noise
vulnerability. Higher-order constellations transmit more information in each
symbol but have smaller decision regions. The adaptive receiver must
therefore select the correct constellation before demodulation.

\section{Noise and Fading}
For AWGN, the received sample is the transmitted sample plus a complex
Gaussian random variable:
\begin{equation}
w[n] \sim \mathcal{CN}(0,\sigma_w^2).
\end{equation}
The noise power is selected from the desired SNR:
\begin{equation}
\sigma_w^2 = \frac{P_s}{10^{\mathrm{SNR}_{dB}/10}},
\end{equation}
where $P_s$ is the received signal power.

Rayleigh fading is appropriate when a dominant line-of-sight path is absent.
The block coefficient used in the simulation is generated as
\begin{equation}
h = \frac{h_I+jh_Q}{\sqrt{2}},
\quad h_I,h_Q \sim \mathcal{N}(0,1).
\end{equation}
The receiver applies the known frame-level channel estimate for equalization:
\begin{equation}
\hat{x}[n] = \frac{y[n]}{h}.
\end{equation}

\section{Carrier Phase and Frequency Offset}
Carrier mismatch rotates the received constellation. In the simulation, the
combined impairment is represented by
\begin{equation}
c[n] = e^{j(\phi+2\pi\Delta f n)}.
\end{equation}
The received signal is corrected using the pilot-assisted estimate available
to the simulation receiver. This design allows the experiment to focus on the
new AMC-to-demodulator control mechanism while preserving a realistic
communication impairment model.

\section{Evaluation Metrics}
\subsection{Bit Error Rate}
Bit error rate measures the fraction of incorrectly recovered bits:
\begin{equation}
\mathrm{BER} = \frac{N_{\mathrm{bit,error}}}{N_{\mathrm{bit,total}}}.
\end{equation}

\subsection{Symbol Error Rate}
Symbol error rate measures the fraction of incorrect constellation decisions:
\begin{equation}
\mathrm{SER} = \frac{N_{\mathrm{symbol,error}}}{N_{\mathrm{symbol,total}}}.
\end{equation}

\subsection{Packet Success and Throughput}
A packet is counted as successful only when all bits in the frame are
recovered correctly. The packet success rate is
\begin{equation}
P_{\mathrm{success}} =
\frac{N_{\mathrm{successful\ packets}}}{N_{\mathrm{transmitted\ packets}}}.
\end{equation}
The reported effective throughput is the average modulation order weighted by
packet success:
\begin{equation}
T_{\mathrm{eff}} =
\left(\frac{1}{M}\sum_{m=1}^{M}b_m\right)P_{\mathrm{success}},
\end{equation}
where $b_m$ is the number of bits per symbol for modulation $m$.

\section{Phase 1 and Phase 2 Distinction}
Phase 1 evaluates whether a CNN can classify a labeled I/Q window. Phase 2
evaluates whether modulation information can be used in a digital link. The
first is a recognition problem; the second is a physical-layer receiver
problem. Consequently, Phase 2 includes transmitter generation, channel
impairments, synchronization, equalization, adaptive demodulation, and
information-recovery metrics.
